This paper examines who benefits more from standardized versus non-standardized tests in university admissions. I focus on a unique event in the UK during the pandemic, when teachers evaluations replaced the final grades of standardized exams, to investigate changes in grade assignment patterns across students and schools. The results show that this method change led to substantial grade inflation with unmasked heterogeneity across students and school attributes. Independent schools were 50\% more likely to assign top grades to their students than public schools. Students taking non-facilitating subjects were 9 percentage point more likely to receive a top grade than students taking facilitating subjects. These findings suggest that incorporating teachers' predictions into college admissions processes rewards students unevenly.